\section{Introduction}

% Paragraph 1: Task and Application
Explaining complex topics in simple terms---the essence of the ELI5 (Explain Like I'm 5) task---is fundamental to education, science communication, and knowledge accessibility. This task targets at generating simplified explanations from complex questions, requiring models to balance factual accuracy with accessibility for non-expert audiences. While large language models have demonstrated impressive capabilities in generating coherent explanations, the task presents a unique challenge: maintaining factual accuracy while achieving simplicity suitable for non-expert audiences. Current approaches predominantly rely on single-pass prompting, where models must simultaneously handle fact retrieval, logical reasoning, and linguistic simplification in one generation step.

% Paragraph 2: Limitations and Root Issue
Single-pass prompting fails to meet the dual requirements of accuracy and simplicity due to three fundamental limitations. First, smaller models often struggle with complex questions due to limited parametric knowledge, leading to incomplete or inaccurate explanations. Second, without external grounding, models may produce hallucinations when facts exceed their training data. Third, single-pass generation offers no explicit decomposition of the reasoning process, making it difficult to diagnose errors or ensure systematic coverage of key concepts. The root technical issue is that monolithic generation cannot simultaneously optimize for factual accuracy, logical reasoning, and linguistic simplification.

% Paragraph 3: Our Technical Solution
We address these challenges through a multi-agent architecture that decomposes explanation generation into five specialized agents (Figure~\ref{fig:pipeline}). First, a \textbf{breakdown agent} decomposes each question into 3--5 targeted search queries for fact retrieval and 3--5 reasoning points for logical analysis, providing a structured roadmap for downstream agents. Second, two agents operate in parallel: a \textbf{reasoning agent} that performs logical analysis of the reasoning points, generating 3--6 logical pathways with conclusions; and a \textbf{scientific agent} that employs hybrid retrieval-augmented generation (RAG) combining BM25 and semantic search over the OpenThoughts-114k dataset, augmented with Wikipedia lookups via parallel requests, extracting 12 grounded facts with citations. Third, a \textbf{synthesis agent} evaluates the quality of both fact extraction and logical reasoning, selecting an adaptive content strategy---reasoning-heavy (70\% reasoning, 30\% facts) when retrieval is weak, facts-heavy (70\% facts, 30\% reasoning) when retrieval is strong, or balanced (50--50) otherwise---and curating 4--6 key points ordered logically. Finally, a \textbf{creative agent} transforms these curated points into accessible ELI5 explanations using higher temperature for natural variation and applying simplification guidelines including everyday vocabulary, relatable analogies, and story-like narrative flow. This staged approach enables explicit reasoning, external grounding, and quality-aware content mixing.

% Paragraph 4: Key Innovation and Benefits
A key innovation is staged batching: by processing all questions together at each agent stage, we reduce vLLM inference calls by approximately 1,000$\times$ compared to per-question processing (Figure~\ref{fig:pipeline}). Processing 1,000 questions requires 4 batched calls instead of 4,000 individual calls, dramatically reducing both latency and cost. The architecture maintains 100\% success rate with acceptable latency (29.64 seconds average) while enabling explicit reasoning decomposition and quality-aware content mixing.

% Paragraph 5: Experimental Summary
To evaluate this architecture, we conduct a comprehensive comparison against baseline single-pass prompting across seven language models spanning 1B to 7B parameters (LLaMA, Qwen, Gemma, Mistral). Using 30,000 samples and eleven evaluation metrics---including LLM-based judges, ROUGE, BERT-Score, and semantic similarity---we uncover a surprising pattern. The multi-agent architecture consistently achieves lower LLM-judged accuracy (-38.2\% on average) while simultaneously improving automatic text quality metrics (+34.2\% ROUGE1, +10.0\% BERT-F1). This accuracy-quality trade-off appears across all tested models, suggesting it is architectural rather than model-specific. Critically, these ROUGE gains are not artifacts of output length: multi-agent explanations are 67.6\% shorter on average while achieving higher lexical overlap with references (Figure~\ref{fig:wordcount}).

% Paragraph 6: Contributions
Our contributions are threefold. First, we introduce a novel five-agent architecture for ELI5 generation featuring adaptive synthesis and efficient staged batching (Section~\ref{sec:methodology}). Second, we provide comprehensive empirical evaluation across diverse models (14 configurations, ~66,000 total samples) and metrics, establishing benchmark comparisons for future work (Section~\ref{sec:experiments}). Third, we discover and analyze the accuracy-quality paradox in multi-agent explanation systems, offering insights into when each architectural approach is appropriate and highlighting important questions about evaluation metric alignment (Section~\ref{sec:discussion}). Our findings suggest that the choice between single-pass and multi-agent approaches should be guided by specific use case priorities rather than universal superiority of either method.
